\answer{ex:18:1}{(а)~$0.9949\le w\le1.0049$, то есть $|1-w|\lesssim0.005$. (б)~$K\ge400$ синапсов.}
\answer{ex:18:2}{$(\omega_R\tau_w)^2=\sqrt{\alpha(\alpha+2+2\varrho)}/\varrho-1$, $\varrho=\tau_m/\tau_w$; $f_R\approx11.1$~Гц; $|Z(\omega_R)|/|Z(0)|=4.6$.}
\answer{ex:18:3}{(а)~$f=1/\bigl(\tau\ln\frac{\mu}{\mu-\theta}\bigr)$; для $1$~Гц нужно $\mu/\theta-1\approx3.7\cdot10^{-44}$. (б)~$T=\pi\tau/\sqrt\mu$, $f\approx3.2$~Гц: $f\propto\sqrt\mu$.}
\answer{ex:18:4}{Интегратор работает при $f\lesssim17.7$~Гц (фильтр нижних частот), резонатор --- только в полосе $5.5$--$22$~Гц (полосовой фильтр).}
\answer{ex:18:5}{(а)~$T_{\min}=\frac{\tau}{w-1}\ln\frac{0.5}{0.3}\approx10.2\ms$. (б)~$B>w-1-0.2=0.3$.}
\answer{ex:18:6}{Постоянная времени подстройки $\tau/(\eta\langle r^2\rangle)$, $\eta=\tau\ln10/60\,\text{с}\approx3.8\cdot10^{-3}$; при $I\neq0$ вес смещается, из $e$ нужно вычесть копию команды $I/\tau$.}
\answer{ex:18:7}{Открытая задача. Резонатор выигрывает лишь в $1.35$ раза при $11$~Гц и проигрывает в $17$ раз при $1$~Гц; основной выигрыш от перестройки --- в пороговом выборе полосы (задача~\ref{ex:18:4}).}
